% =============================================================================
% DYN-AXIOME (DYNAMIK) - Vollständige Ausformulierung
% BCM 2.2 META-Modul
% =============================================================================

\section{Dynamik-Axiome (MA-DYN)}

Die DYN-Axiome definieren zeitliche Prozesse, Veränderungen und Evolution.
Sie beantworten die Frage: \textit{Wie verändert es sich?}

%==============================================================================
% MA-DYN-01: Zeit ist kontinuierlich
%==============================================================================

\subsubsection*{MA-DYN-01: Zeit ist kontinuierlich}

\textbf{Statement:} Verhalten findet in kontinuierlicher Zeit statt.

\textbf{Formel:}
\begin{equation}
\tau \in \mathbb{R}^+
\end{equation}

\textbf{Beschreibung:}
Zeit ist nicht diskret, sondern fließend:
\begin{itemize}
    \item Zustände ändern sich stetig
    \item Übergänge sind graduell
    \item Differentialgleichungen sind angemessene Modelle
\end{itemize}

\textbf{BCM-Relevanz:}
Ermöglicht dynamische Modellierung mit kontinuierlichen Funktionen. Habituation, Lernen und Evolution sind als stetige Prozesse modellierbar.

\textbf{BEATRIX-Relevanz:}
BEATRIX arbeitet mit diskreten Zeitpunkten (Snapshot), aber:
\begin{itemize}
    \item Interpolation zwischen Zeitpunkten möglich
    \item Trajektorien statt Zustände modellieren
    \item Zeitliche Glättung bei Messungen
\end{itemize}

\textbf{Konsequenz bei Fehlen:}
Nur diskrete Zustände. Keine kontinuierliche Dynamik. Sprunghafte Übergänge statt gradueller.

\textbf{Validiert:} WTX

\textbf{Abhängigkeiten:} MA-ONT-05

\textbf{Literatur:} Newton (1687); Strogatz (2015)

%==============================================================================
% MA-DYN-02: Diskrete Events
%==============================================================================

\subsubsection*{MA-DYN-02: Diskrete Events}

\textbf{Statement:} Entscheidungen sind diskrete Ereignisse in kontinuierlicher Zeit.

\textbf{Formel:}
\begin{equation}
\text{Decision}_i \text{ at } \tau_i \in T
\end{equation}

\textbf{Beschreibung:}
Obwohl Zeit kontinuierlich ist, sind Entscheidungen punktuell:
\begin{itemize}
    \item Kaufentscheidung: Zeitpunkt $\tau_k$
    \item Wahlentscheidung: Zeitpunkt $\tau_w$
    \item Sequenz von Decisions: $\{\tau_1, \tau_2, ..., \tau_n\}$
\end{itemize}

\textbf{BCM-Relevanz:}
Verbindet kontinuierliche Dynamik mit diskreten Handlungen. Interventionen zielen auf Entscheidungszeitpunkte.

\textbf{BEATRIX-Relevanz:}
BEATRIX fokussiert auf Decision Points:
\begin{itemize}
    \item Wann fällt die Entscheidung?
    \item Interventions-Timing optimieren
    \item Sequenz von Entscheidungen modellieren
\end{itemize}

\textbf{Konsequenz bei Fehlen:}
Keine Unterscheidung Prozess/Event. Entscheidungen wären stetig statt punktuell.

\textbf{Validiert:} WTX

\textbf{Abhängigkeiten:} MA-DYN-01

\textbf{Literatur:} Ross (1983)

%==============================================================================
% MA-DYN-03: Pfadabhängigkeit
%==============================================================================

\subsubsection*{MA-DYN-03: Pfadabhängigkeit}

\textbf{Statement:} Geschichte beeinflusst Gegenwart.

\textbf{Formel:}
\begin{equation}
\text{State}(\tau) = f(\text{History}(0..\tau))
\end{equation}

\textbf{Beschreibung:}
Der aktuelle Zustand ist nicht nur von aktuellen Bedingungen abhängig:
\begin{itemize}
    \item Frühere Entscheidungen schaffen Pfade
    \item Lock-In-Effekte: QWERTY-Tastatur
    \item Gewohnheiten sind akkumulierte Geschichte
\end{itemize}

\textbf{BCM-Relevanz:}
Erklärt Persistenz von Verhalten. Verhaltensänderung erfordert ``Pfadbruch''. Frühe Intervention hat höheren Hebel.

\textbf{BEATRIX-Relevanz:}
BEATRIX berücksichtigt Verhaltensgeschichte:
\begin{itemize}
    \item Wie lange besteht das Verhalten schon?
    \item Welche Lock-Ins existieren?
    \item Historische Daten für bessere Vorhersage nutzen
\end{itemize}

\textbf{Konsequenz bei Fehlen:}
Zustand nur von aktuellen Bedingungen abhängig. Markov-Annahme. Geschichte irrelevant.

\textbf{Validiert:} WTX

\textbf{Abhängigkeiten:} MA-DYN-01

\textbf{Literatur:} Arthur (1989); David (1985)

%==============================================================================
% MA-DYN-04: Multiple Equilibria
%==============================================================================

\subsubsection*{MA-DYN-04: Multiple Equilibria}

\textbf{Statement:} Es kann mehrere stabile Zustände geben.

\textbf{Formel:}
\begin{equation}
|\{\text{State}: \frac{d\text{State}}{dt} = 0\}| \geq 1
\end{equation}

\textbf{Beschreibung:}
Systeme können mehrere Gleichgewichte haben:
\begin{itemize}
    \item Kooperation vs. Defektion als stabile Zustände
    \item ``Gute'' vs. ``schlechte'' Gewohnheit
    \item Gesellschaftliche Normen A vs. B
\end{itemize}
Welches Gleichgewicht erreicht wird, hängt von Anfangsbedingungen ab.

\textbf{BCM-Relevanz:}
Erklärt, warum ähnliche Populationen unterschiedliche Normen haben können. Interventionen können Gleichgewichtswechsel auslösen.

\textbf{BEATRIX-Relevanz:}
BEATRIX analysiert Gleichgewichte:
\begin{itemize}
    \item In welchem Gleichgewicht befindet sich die Persona/Gruppe?
    \item Ist ein Wechsel zu einem besseren Gleichgewicht möglich?
    \item Welche Intervention löst den Wechsel aus?
\end{itemize}

\textbf{Konsequenz bei Fehlen:}
Einziges Gleichgewicht. Keine alternativen Attraktoren. Deterministische Konvergenz.

\textbf{Validiert:} WTX

\textbf{Abhängigkeiten:} MA-DYN-03

\textbf{Literatur:} Schelling (1978); Young (1998)

%==============================================================================
% MA-DYN-05: Hysterese
%==============================================================================

\subsubsection*{MA-DYN-05: Hysterese möglich}

\textbf{Statement:} Systemzustand hängt von der Richtung der Veränderung ab.

\textbf{Formel:}
\begin{equation}
\text{State}(\uparrow) \neq \text{State}(\downarrow) \quad \text{bei gleichem Trigger}
\end{equation}

\textbf{Beschreibung:}
Der Weg zählt:
\begin{itemize}
    \item Gewohnheit aufbauen ist leichter als abbauen
    \item Vertrauen aufbauen dauert, zerstören geht schnell
    \item Norm-Etablierung vs. Norm-Erosion asymmetrisch
\end{itemize}

\textbf{BCM-Relevanz:}
Asymmetrie von Aufbau und Abbau. Positive Verhaltensänderung kann andere Schwellen haben als Rückfall.

\textbf{BEATRIX-Relevanz:}
BEATRIX modelliert Hysterese:
\begin{itemize}
    \item Unterschiedliche Schwellen für Veränderung vs. Rückfall
    \item Stabilisierung nach Verhaltensänderung nötig
    \item Warnung vor ``zu schnellem'' Absetzen von Interventionen
\end{itemize}

\textbf{Konsequenz bei Fehlen:}
Symmetrische Dynamik angenommen. Richtungsunabhängigkeit.

\textbf{Validiert:} WTX

\textbf{Abhängigkeiten:} MA-DYN-03

\textbf{Literatur:} Scheffer (2009)

%==============================================================================
% MA-DYN-06: Tipping Points
%==============================================================================

\subsubsection*{MA-DYN-06: Tipping Points existieren}

\textbf{Statement:} Es gibt kritische Schwellen für qualitative Zustandswechsel.

\textbf{Formel:}
\begin{equation}
\exists \theta: \text{State}(x < \theta) \neq \text{State}(x > \theta) \quad \text{qualitativ}
\end{equation}

\textbf{Beschreibung:}
Kleine Änderungen können große Wirkung haben -- wenn sie eine Schwelle überschreiten:
\begin{itemize}
    \item Epidemiologischer Tipping Point: $R > 1$
    \item Soziale Tipping Points: Norm-Kaskaden
    \item Individuelle Tipping Points: Verhaltensumbruch
\end{itemize}

\textbf{BCM-Relevanz:}
Erklärt Nicht-Linearität. Kleine Interventionen können große Effekte haben -- wenn sie den Tipping Point erreichen.

\textbf{BEATRIX-Relevanz:}
BEATRIX identifiziert Tipping Points:
\begin{itemize}
    \item Wie nah ist das System am Tipping Point?
    \item Welche Intervention erreicht den Kipppunkt?
    \item Warnung vor instabilen Zuständen
\end{itemize}

\textbf{Konsequenz bei Fehlen:}
Lineare Extrapolation. Sprunghafte Veränderungen unerwartet.

\textbf{Validiert:} WTX

\textbf{Abhängigkeiten:} MA-DYN-04

\textbf{Literatur:} Gladwell (2000); Scheffer et al. (2012)

%==============================================================================
% MA-DYN-07: Gewohnheitsformation
%==============================================================================

\subsubsection*{MA-DYN-07: Gewohnheitsformation}

\textbf{Statement:} Wiederholung führt zu Automatisierung.

\textbf{Formel:}
\begin{equation}
P(V_{\text{auto}}|\text{Repetition}) \uparrow \text{ mit } n
\end{equation}

\textbf{Beschreibung:}
Gewohnheiten bilden sich durch Wiederholung:
\begin{itemize}
    \item Ca. 66 Tage für automatisches Verhalten (Lally et al.)
    \item Exponentieller Anstieg der Automatisierung
    \item Plateau nach Stabilisierung
\end{itemize}

\textbf{BCM-Relevanz:}
Ziel vieler Interventionen: Gewohnheitsbildung. Einmalige Verhaltensänderung ist weniger wertvoll als dauerhafte.

\textbf{BEATRIX-Relevanz:}
BEATRIX plant für Gewohnheitsbildung:
\begin{itemize}
    \item Wie oft muss das Verhalten wiederholt werden?
    \item Welche Cues etablieren die Gewohnheit?
    \item Tracking der Gewohnheitsstärke
\end{itemize}

\textbf{Konsequenz bei Fehlen:}
Keine Automatisierung. Jede Entscheidung deliberativ.

\textbf{Validiert:} WTX

\textbf{Abhängigkeiten:} MA-EPI-21

\textbf{Literatur:} Lally et al. (2010); Wood \& Rünger (2016)

%==============================================================================
% MA-DYN-08: Habit Loop
%==============================================================================

\subsubsection*{MA-DYN-08: Habit Loop}

\textbf{Statement:} Gewohnheiten folgen dem Cue-Routine-Reward-Muster.

\textbf{Formel:}
\begin{equation}
\text{Habit} = f(\text{Cue}, \text{Routine}, \text{Reward})
\end{equation}

\textbf{Beschreibung:}
Der Habit Loop nach Duhigg:
\begin{enumerate}
    \item \textbf{Cue}: Trigger (Zeit, Ort, Emotion, andere Personen)
    \item \textbf{Routine}: Das automatische Verhalten
    \item \textbf{Reward}: Belohnung, die den Loop verstärkt
\end{enumerate}

\textbf{BCM-Relevanz:}
Interventionsansatz: Bestehenden Cue mit neuer Routine verbinden. ``Habit stacking''. Reward engineering.

\textbf{BEATRIX-Relevanz:}
BEATRIX analysiert Habit Loops:
\begin{itemize}
    \item Welche Cues triggern welches Verhalten?
    \item Welche Rewards verstärken es?
    \item Wie kann der Loop umgestaltet werden?
\end{itemize}

\textbf{Konsequenz bei Fehlen:}
Keine Struktur für Gewohnheitsänderung. Nur globale Verhaltensänderung statt gezielter Manipulation.

\textbf{Validiert:} WTX

\textbf{Abhängigkeiten:} MA-DYN-07

\textbf{Literatur:} Duhigg (2012); Wood \& Neal (2007)

%==============================================================================
% MA-DYN-09: Commitment Devices
%==============================================================================

\subsubsection*{MA-DYN-09: Commitment Devices}

\textbf{Statement:} Selbstbindung ist möglich und kann rational sein.

\textbf{Formel:}
\begin{equation}
U(\text{Commit}) > U(\text{NoCommit}) \quad \text{wenn } \delta_{\text{now}} > \delta_{\text{future}}
\end{equation}

\textbf{Beschreibung:}
Menschen können sich selbst binden, um zukünftige Schwäche zu überwinden:
\begin{itemize}
    \item Ulysses und die Sirenen
    \item Sparverträge mit Strafe bei vorzeitiger Auflösung
    \item Öffentliche Ankündigungen
\end{itemize}

\textbf{BCM-Relevanz:}
Lösung für zeitinkonsistente Präferenzen (MA-DYN-14). Commitment ist keine Irrationalität, sondern Metarationalität.

\textbf{BEATRIX-Relevanz:}
BEATRIX empfiehlt Commitment Devices:
\begin{itemize}
    \item Welche Form der Selbstbindung passt zur Persona?
    \item Soziale vs. finanzielle Commitments
    \item Stärke des Commitments kalibrieren
\end{itemize}

\textbf{Konsequenz bei Fehlen:}
Keine Lösung für Present Bias. Zeitinkonsistenz bleibt unaddressiert.

\textbf{Validiert:} WAX

\textbf{Abhängigkeiten:} MA-EPI-16

\textbf{Literatur:} Elster (1979); Bryan et al. (2010)

%==============================================================================
% MA-DYN-10: Soziale Diffusion
%==============================================================================

\subsubsection*{MA-DYN-10: Soziale Diffusion}

\textbf{Statement:} Verhalten breitet sich sozial aus.

\textbf{Formel:}
\begin{equation}
P(V_i|V_{\text{neighbors}}) = f(\text{Density}(V_{\text{neighbors}}))
\end{equation}

\textbf{Beschreibung:}
Verhalten ist ansteckend:
\begin{itemize}
    \item Peer Effects: Freunde beeinflussen Verhalten
    \item Netzwerk-Diffusion: Von Early Adopters zur Mehrheit
    \item Kaskaden: Kleine Anfänge, große Wirkung
\end{itemize}

\textbf{BCM-Relevanz:}
Grundlage für SOCIAL-Interventionen. Der $\lambda$-Parameter hängt von wahrgenommener Kooperation anderer ab.

\textbf{BEATRIX-Relevanz:}
BEATRIX modelliert Netzwerkeffekte:
\begin{itemize}
    \item Wer sind die relevanten Peers der Persona?
    \item Welche Diffusionsstrategien sind effektiv?
    \item Seed-Auswahl für maximale Verbreitung
\end{itemize}

\textbf{Konsequenz bei Fehlen:}
Individualistisches Modell. Soziale Ansteckung nicht abbildbar.

\textbf{Validiert:} KNU

\textbf{Abhängigkeiten:} MA-ONT-10

\textbf{Literatur:} Rogers (2003); Centola \& Macy (2007)

%==============================================================================
% MA-DYN-11: Bereits in axiom_beispiele.tex
%==============================================================================

%==============================================================================
% MA-DYN-12: Norm-Kaskaden
%==============================================================================

\subsubsection*{MA-DYN-12: Norm-Kaskaden}

\textbf{Statement:} Normen können rapid kippen.

\textbf{Formel:}
\begin{equation}
\frac{dN}{dt} \text{ kann sprunghaft sein}
\end{equation}

\textbf{Beschreibung:}
Normenwandel ist oft nicht graduell:
\begin{itemize}
    \item Raucher-Norm: von ``cool'' zu ``uncool'' in einer Generation
    \item Gleichgeschlechtliche Ehe: rapide Akzeptanzsteigerung
    \item ``Preference falsification'' → plötzliches Offenlegen
\end{itemize}

\textbf{BCM-Relevanz:}
Dynamische Ergänzung zu MA-ONT-24 (Normen evolvieren). Nicht nur dass, sondern wie: Sprünge statt lineare Trends.

\textbf{BEATRIX-Relevanz:}
BEATRIX erkennt Kaskadenrisiken:
\begin{itemize}
    \item Ist eine Norm-Kaskade im Gange?
    \item Kann eine Intervention die Kaskade auslösen?
    \item Timing für maximalen Impact
\end{itemize}

\textbf{Konsequenz bei Fehlen:}
Nur gradueller Normenwandel modelliert. Plötzliche Umbrüche überraschen.

\textbf{Validiert:} KNU

\textbf{Abhängigkeiten:} MA-ONT-24, MA-DYN-11

\textbf{Literatur:} Kuran (1995); Bicchieri (2016)

%==============================================================================
% MA-DYN-13: Discounting
%==============================================================================

\subsubsection*{MA-DYN-13: Discounting}

\textbf{Statement:} Zukünftiger Nutzen wird diskontiert.

\textbf{Formel:}
\begin{equation}
U_{\text{present}}(U_{\text{future}}) = \delta^t \cdot U_{\text{future}}
\end{equation}

\textbf{Beschreibung:}
Ein Euro heute ist mehr wert als ein Euro morgen:
\begin{itemize}
    \item Zeitpräferenz: Gegenwart wird vorgezogen
    \item Unsicherheit: Zukunft ist ungewiss
    \item Opportunity Cost: Verzicht auf heutige Nutzung
\end{itemize}

\textbf{BCM-Relevanz:}
Fundamental für intertemporale Entscheidungen. Erklärt Unterinvestition in Zukunft (Sparen, Prävention, Klimaschutz).

\textbf{BEATRIX-Relevanz:}
BEATRIX schätzt Diskontrate:
\begin{itemize}
    \item Wie stark diskontiert die Persona?
    \item Kurzfristige Belohnungen vs. langfristiger Nutzen
    \item Interventionen mit optimiertem zeitlichen Profil
\end{itemize}

\textbf{Konsequenz bei Fehlen:}
Keine Zeitpräferenz. Gegenwart und Zukunft gleichwertig.

\textbf{Validiert:} INU

\textbf{Abhängigkeiten:} MA-EPI-16

\textbf{Literatur:} Samuelson (1937); Frederick et al. (2002)

%==============================================================================
% MA-DYN-14: Bereits in axiom_beispiele.tex
%==============================================================================

%==============================================================================
% MA-DYN-15: Projection Bias
%==============================================================================

\subsubsection*{MA-DYN-15: Projection Bias}

\textbf{Statement:} Aktuelle Zustände werden in die Zukunft projiziert.

\textbf{Formel:}
\begin{equation}
E[U_{\text{future}}|\text{State}_{\text{now}}] \text{ biased toward State}_{\text{now}}
\end{equation}

\textbf{Beschreibung:}
Menschen überschätzen die Persistenz aktueller Zustände:
\begin{itemize}
    \item Hungrig einkaufen: Zu viel kaufen
    \item Gute Laune: Optimistische Prognosen
    \item Schlechte Laune: Pessimistische Prognosen
\end{itemize}

\textbf{BCM-Relevanz:}
Erklärt systematische Fehlplanung. Zukünftige Zustände (Hunger, Motivation, Energie) werden falsch antizipiert.

\textbf{BEATRIX-Relevanz:}
BEATRIX korrigiert für Projection Bias:
\begin{itemize}
    \item In welchem Zustand befindet sich die Persona jetzt?
    \item Wie würde sie in anderem Zustand entscheiden?
    \item Interventionen nicht in extremen Zuständen planen
\end{itemize}

\textbf{Konsequenz bei Fehlen:}
Aktuelle Zustände würden als repräsentativ für zukünftige gelten.

\textbf{Validiert:} INU

\textbf{Abhängigkeiten:} MA-EPI-06

\textbf{Literatur:} Loewenstein et al. (2003)

%==============================================================================
% MA-DYN-16: Hot-Cold Empathy Gap
%==============================================================================

\subsubsection*{MA-DYN-16: Hot-Cold Empathy Gap}

\textbf{Statement:} In ``cold'' Zuständen werden ``hot'' Zustände unterschätzt.

\textbf{Formel:}
\begin{equation}
E[U_{\text{hot}}|\text{State}_{\text{cold}}] < U_{\text{hot, actual}}
\end{equation}

\textbf{Beschreibung:}
Die Intensität emotionaler Zustände wird nicht antizipiert:
\begin{itemize}
    \item Nüchtern: Unterschätzen von Alkohol-Verlangen
    \item Satt: Unterschätzen von Hunger
    \item Ruhig: Unterschätzen von Wut
\end{itemize}

\textbf{BCM-Relevanz:}
Erklärt Scheitern von ``cold'' geplanten Vorsätzen in ``hot'' Situationen. Commitment Devices helfen, wenn Cold-Self den Hot-Self bindet.

\textbf{BEATRIX-Relevanz:}
BEATRIX berücksichtigt Hot-Cold Gap:
\begin{itemize}
    \item Plane Interventionen für den ``hot'' Zustand
    \item Commitments in ``cold'' Zuständen eingehen lassen
    \item Warnung vor Hot-State-Entscheidungen
\end{itemize}

\textbf{Konsequenz bei Fehlen:}
Emotionale Zustände würden korrekt antizipiert. Vorsätze würden immer halten.

\textbf{Validiert:} INU, AWX

\textbf{Abhängigkeiten:} MA-DYN-15

\textbf{Literatur:} Loewenstein (1996, 2005)

%==============================================================================
% MA-DYN-17: Ego Depletion
%==============================================================================

\subsubsection*{MA-DYN-17: Ego Depletion}

\textbf{Statement:} Selbstkontrolle ist eine erschöpfliche Ressource.

\textbf{Formel:}
\begin{equation}
\text{Willpower}(t) \downarrow \text{ mit Nutzung}
\end{equation}

\textbf{Beschreibung:}
Selbstkontrolle funktioniert wie ein Muskel:
\begin{itemize}
    \item Ermüdet durch Nutzung
    \item Regeneriert durch Ruhe
    \item Kann durch Training gestärkt werden
\end{itemize}
Hinweis: Effekt ist umstritten (Replication Crisis), aber praktisch relevant.

\textbf{BCM-Relevanz:}
Erklärt Selbstkontrollversagen am Abend, nach Stress. Interventionen sollten Willpower-Status berücksichtigen.

\textbf{BEATRIX-Relevanz:}
BEATRIX berücksichtigt Ego Depletion:
\begin{itemize}
    \item Wann ist die Persona wahrscheinlich depleted?
    \item Interventionen mit geringem Willpower-Bedarf bevorzugen
    \item Timing für Willpower-intensive Entscheidungen
\end{itemize}

\textbf{Konsequenz bei Fehlen:}
Konstante Selbstkontrolle angenommen. Ermüdung ignoriert.

\textbf{Validiert:} WAX, WTX

\textbf{Abhängigkeiten:} MA-EPI-05

\textbf{Literatur:} Baumeister et al. (1998); (umstritten: Carter \& McCullough, 2014)

%==============================================================================
% MA-DYN-18: Implementation Intentions
%==============================================================================

\subsubsection*{MA-DYN-18: Implementation Intentions}

\textbf{Statement:} Wenn-Dann-Pläne erhöhen Umsetzungswahrscheinlichkeit.

\textbf{Formel:}
\begin{equation}
P(V|\text{If-Then Plan}) > P(V|\text{Goal only})
\end{equation}

\textbf{Beschreibung:}
Konkrete Wenn-Dann-Pläne schließen die Intention-Action-Gap:
\begin{itemize}
    \item ``Wenn X, dann tue ich Y''
    \item Automatische Aktivierung durch Cue
    \item Reduziert deliberativen Aufwand
\end{itemize}

\textbf{BCM-Relevanz:}
Effektive Intervention für WTX-Defizite. Kombiniert gut mit NUDGE (Default als Implementation Intention).

\textbf{BEATRIX-Relevanz:}
BEATRIX empfiehlt Implementation Intentions:
\begin{itemize}
    \item Konkrete Wenn-Dann-Formulierungen generieren
    \item Passende Cues für die Persona identifizieren
    \item Als Teil von BOOST-Interventionen
\end{itemize}

\textbf{Konsequenz bei Fehlen:}
Nur Ziele ohne Umsetzungspläne. Intention-Action-Gap bleibt.

\textbf{Validiert:} WTX

\textbf{Abhängigkeiten:} MA-EPI-11

\textbf{Literatur:} Gollwitzer (1999); Gollwitzer \& Sheeran (2006)

%==============================================================================
% MA-DYN-19: Adaptation
%==============================================================================

\subsubsection*{MA-DYN-19: Adaptation}

\textbf{Statement:} Nutzen adaptiert an neuen Status Quo.

\textbf{Formel:}
\begin{equation}
U(\text{State}) \rightarrow U_{\text{baseline}} \quad \text{as } \tau \rightarrow \infty
\end{equation}

\textbf{Beschreibung:}
Menschen gewöhnen sich an alles:
\begin{itemize}
    \item Positive Veränderungen: Freude verblasst
    \item Negative Veränderungen: Leid verblasst
    \item Neuer Normalzustand wird Referenzpunkt
\end{itemize}

\textbf{BCM-Relevanz:}
Erklärt das Hedonic Treadmill (MA-DYN-20). Dauerhafte Nutzenänderung ist schwierig zu erreichen.

\textbf{BEATRIX-Relevanz:}
BEATRIX berücksichtigt Adaptation:
\begin{itemize}
    \item Erwartete Adaptationsrate modellieren
    \item Nachhaltige vs. temporäre Nutzeneffekte
    \item Intervention gegen Adaptation designen
\end{itemize}

\textbf{Konsequenz bei Fehlen:}
Permanente Nutzenveränderungen angenommen. Überschätzung langfristiger Effekte.

\textbf{Validiert:} INU

\textbf{Abhängigkeiten:} MA-ONT-20

\textbf{Literatur:} Helson (1964); Wilson \& Gilbert (2005)

%==============================================================================
% MA-DYN-20: Hedonic Treadmill
%==============================================================================

\subsubsection*{MA-DYN-20: Hedonic Treadmill}

\textbf{Statement:} Glück kehrt zu einer individuellen Baseline zurück.

\textbf{Formel:}
\begin{equation}
\text{Happiness} \rightarrow \text{Baseline} \quad \text{nach positiven/negativen Events}
\end{equation}

\textbf{Beschreibung:}
Das hedonistische Laufband:
\begin{itemize}
    \item Lottogewinner: Nach 1 Jahr wieder auf Ausgangsniveau
    \item Querschnittslähmung: Nach 2 Jahren signifikante Erholung
    \item Set Point Theory: Genetisch determinierte Baseline
\end{itemize}

\textbf{BCM-Relevanz:}
Relativiert materielle Interventionen. Mehr Geld macht nicht dauerhaft glücklicher. Fokus auf prozessuale statt ergebnisbasierte Interventionen.

\textbf{BEATRIX-Relevanz:}
BEATRIX berücksichtigt Hedonic Adaptation:
\begin{itemize}
    \item Welche Interventionen haben nachhaltige Effekte?
    \item Prozess-Glück vs. Ergebnis-Glück
    \item Warnung vor Überschätzung materieller Anreize
\end{itemize}

\textbf{Konsequenz bei Fehlen:}
Lineare Beziehung Ressourcen → Glück. Materielle Interventionen überschätzt.

\textbf{Validiert:} INU

\textbf{Abhängigkeiten:} MA-DYN-19

\textbf{Literatur:} Brickman \& Campbell (1971); Diener et al. (2006)

%==============================================================================
% MA-DYN-21: Peak-End Rule
%==============================================================================

\subsubsection*{MA-DYN-21: Peak-End Rule}

\textbf{Statement:} Erinnerung = Peak + End.

\textbf{Formel:}
\begin{equation}
\text{Recalled}_U = f(\max_U, \text{End}_U)
\end{equation}

\textbf{Beschreibung:}
Die Erinnerung an eine Erfahrung wird dominiert von:
\begin{itemize}
    \item Dem intensivsten Moment (Peak)
    \item Dem Endmoment
    \item Dauer spielt kaum eine Rolle
\end{itemize}

\textbf{BCM-Relevanz:}
Erklärt, warum Erfahrungen anders erinnert werden als erlebt. Interventionen können das Ende optimieren für bessere Gesamtbewertung.

\textbf{BEATRIX-Relevanz:}
BEATRIX nutzt Peak-End Rule:
\begin{itemize}
    \item Positive Enden für Interventionserfahrungen designen
    \item Negative Peaks vermeiden
    \item Dauer ist weniger wichtig als Qualität der Momente
\end{itemize}

\textbf{Konsequenz bei Fehlen:}
Dauer-gewichtete Durchschnitte angenommen. Peak und End-Effekte ignoriert.

\textbf{Validiert:} INU

\textbf{Abhängigkeiten:} MA-EPI-12

\textbf{Literatur:} Kahneman et al. (1993); Kahneman (1999)

%==============================================================================
% MA-DYN-22: Identitäts-Shifts
%==============================================================================

\subsubsection*{MA-DYN-22: Identitäts-Shifts}

\textbf{Statement:} Identitäten können sich fundamental ändern.

\textbf{Formel:}
\begin{equation}
I(\tau + \Delta\tau) \not\subset I(\tau) \quad \text{möglich}
\end{equation}

\textbf{Beschreibung:}
Nicht nur Evolution, sondern Transformation:
\begin{itemize}
    \item Konversion: Religiöse, politische
    \item Life Events: Elternschaft, Scheidung, Krankheit
    \item ``Ich bin nicht mehr derselbe Mensch''
\end{itemize}

\textbf{BCM-Relevanz:}
Ermöglicht Modellierung fundamentaler Identitätsänderungen. IDENTITY-Interventionen können auf Shift zielen, nicht nur Salienz.

\textbf{BEATRIX-Relevanz:}
BEATRIX erkennt Identitäts-Shift-Potenzial:
\begin{itemize}
    \item Steht die Persona vor einem Life Event?
    \item Gibt es Anzeichen für Identitätskrise?
    \item Interventionen für Übergangszeiten optimieren
\end{itemize}

\textbf{Konsequenz bei Fehlen:}
Nur graduelle Identitätsänderung. Fundamentale Transformationen nicht modellierbar.

\textbf{Validiert:} IDN

\textbf{Abhängigkeiten:} MA-ONT-23

\textbf{Literatur:} James (1902); Paloutzian (2005)

%==============================================================================
% MA-DYN-23: λ-Dynamik
%==============================================================================

\subsubsection*{MA-DYN-23: $\lambda$-Dynamik}

\textbf{Statement:} Der Kooperationsparameter $\lambda$ verändert sich über Zeit.

\textbf{Formel:}
\begin{equation}
\lambda(\tau + \Delta\tau) = f(\lambda(\tau), \text{Experience}, \text{Context})
\end{equation}

\textbf{Beschreibung:}
$\lambda$ ist nicht statisch:
\begin{itemize}
    \item Positive Kooperationserfahrungen erhöhen $\lambda$
    \item Betrug und Enttäuschung senken $\lambda$
    \item Gesellschaftlicher Wandel beeinflusst Basis-$\lambda$
\end{itemize}

\textbf{BCM-Relevanz:}
Dynamische Erweiterung des $\lambda$-Konzepts. Erklärt, wie Kooperationsbereitschaft sich entwickelt.

\textbf{BEATRIX-Relevanz:}
BEATRIX modelliert $\lambda$-Entwicklung:
\begin{itemize}
    \item Wie hat sich $\lambda$ entwickelt?
    \item Welche Erfahrungen beeinflussen es?
    \item Kann $\lambda$ durch Intervention erhöht werden?
\end{itemize}

\textbf{Konsequenz bei Fehlen:}
Statisches $\lambda$ als Trait. Keine Kooperations-Entwicklung modellierbar.

\textbf{Validiert:} WAX

\textbf{Abhängigkeiten:} MA-ONT-14

\textbf{Literatur:} Ostrom (1998); Fehr \& Gächter (2000)

%==============================================================================
% MA-DYN-24: System-Evolution
%==============================================================================

\subsubsection*{MA-DYN-24: System-Evolution}

\textbf{Statement:} Das BCM-Gesamtsystem evolviert.

\textbf{Formel:}
\begin{equation}
\Phi(\tau + \Delta\tau) = \Phi(\tau) + \mu \cdot \Delta\Phi
\end{equation}

\textbf{Beschreibung:}
Nicht nur Parameter ändern sich, sondern das System als Ganzes:
\begin{itemize}
    \item Neue Verhaltenskategorien entstehen (Digital Behavior)
    \item Neue Interventionstypen werden entdeckt
    \item Systemgrenzen verschieben sich
\end{itemize}

\textbf{BCM-Relevanz:}
Meta-Ebene: Das BCM selbst ist nicht statisch. Neue Forschung erweitert das Modell. Version 2.2 ≠ Version 3.0.

\textbf{BEATRIX-Relevanz:}
BEATRIX wird kontinuierlich aktualisiert:
\begin{itemize}
    \item Neue Erkenntnisse integrieren
    \item Modellparameter rekalibrieren
    \item Systementwicklung dokumentieren
\end{itemize}

\textbf{Konsequenz bei Fehlen:}
Statisches Modell. Keine Weiterentwicklung vorgesehen.

\textbf{Validiert:} META

\textbf{Abhängigkeiten:} MA-ONT-25

\textbf{Literatur:} Kuhn (1962); Lakatos (1978)

%==============================================================================
% MA-DYN-25: Meta-Stabilität
%==============================================================================

\subsubsection*{MA-DYN-25: Meta-Stabilität}

\textbf{Statement:} Fundamentale Struktur bleibt erhalten, auch wenn Inhalte sich ändern.

\textbf{Formel:}
\begin{equation}
\frac{\partial \text{Structure}}{\partial \tau} \rightarrow 0 \quad \text{auch wenn } \frac{\partial \text{Content}}{\partial \tau} \neq 0
\end{equation}

\textbf{Beschreibung:}
Trotz aller Dynamik gibt es Invarianten:
\begin{itemize}
    \item Menschen haben Präferenzen (welche auch immer)
    \item Soziale Normen existieren (welche auch immer)
    \item Biases sind systematisch (welche auch immer)
\end{itemize}
Die Struktur des BCM bleibt stabil, die Inhalte variieren.

\textbf{BCM-Relevanz:}
Ontologisches Fundament. Das BCM ist nicht nur ein historisches Snapshot, sondern beschreibt tiefere Strukturen, die invariant sind.

\textbf{BEATRIX-Relevanz:}
BEATRIX verlässt sich auf Meta-Stabilität:
\begin{itemize}
    \item Strukturannahmen sind robust
    \item Parameter-Aktualisierung, nicht Modell-Austausch
    \item Vertrauen in Grundarchitektur
\end{itemize}

\textbf{Konsequenz bei Fehlen:}
Alles wäre in Fluss. Kein stabiler Rahmen für Modellierung. Beliebigkeit.

\textbf{Validiert:} META

\textbf{Abhängigkeiten:} MA-DYN-24

\textbf{Literatur:} Eigen (1971); Kauffman (1993)

\bigskip
\noindent\rule{\textwidth}{0.4pt}
\textit{Ende der DYN-Axiome. Alle 25 Dynamik-Axiome sind nun vollständig dokumentiert.}
